% =====================================================
% 题型：平面向量单位向量夹角选择题 (q_vector_unit_angle.tex)
% =====================================================
% 难度：★ (基础)
% =====================================================
\newcommand{\qVectorUnitAngleStem}{若单位向量 \(\vec{e}_1, \vec{e}_2\) 满足 \(|\vec{e}_1 + \vec{e}_2| = 1\)，则 \(\vec{e}_1, \vec{e}_2\) 的夹角为 \hfill （\quad）}

\newcommand{\qVectorUnitAngleOptions}{%
  \mcoptions[4]{\(\dfrac{\pi}{6}\)}{\(\dfrac{\pi}{3}\)}{\(\dfrac{\pi}{2}\)}{\(\dfrac{2\pi}{3}\)}%
}

\newcommand{\qVectorUnitAngleAnswer}{D}

\newcommand{\qVectorUnitAngleSolution}{%
  【解析方法一：代数平方化】
  设 \(\vec{e}_1, \vec{e}_2\) 的夹角为 \(\theta\)。因为 \(\vec{e}_1, \vec{e}_2\) 是单位向量，所以 \(|\vec{e}_1| = |\vec{e}_2| = 1\)。

  对已知等式 \(|\vec{e}_1 + \vec{e}_2| = 1\) 两边平方得：
  \[
    |\vec{e}_1 + \vec{e}_2|^2 = |\vec{e}_1|^2 + |\vec{e}_2|^2 + 2\vec{e}_1 \cdot \vec{e}_2 = 1.
  \]
  代入模长值并展开数量积：
  \[
    1^2 + 1^2 + 2 |\vec{e}_1| |\vec{e}_2| \cos\theta = 1 \implies 2 + 2\cos\theta = 1 \implies 2\cos\theta = -1 \implies \cos\theta = -\frac{1}{2}.
  \]
  因为向量夹角 \(\theta \in [0, \pi]\)，所以：
  \[
    \theta = \frac{2\pi}{3}.
  \]
  故选 D。

  【解析方法二：几何向量合成】
  若将两单位向量 \(\vec{e}_1, \vec{e}_2\) 首尾相接，由于它们的模长均为 \(1\)，且合成向量 \(\vec{e}_1 + \vec{e}_2\) 的模长也是 \(1\)，这三个向量正好构成一个边长为 \(1\) 的等边三角形。

  此时合成向量与分向量的夹角为 \(\dfrac{\pi}{3}\)（即三角形的内角）。但注意，两向量 \(\vec{e}_1, \vec{e}_2\) 的夹角是指将它们平移至共同起点时，两射线所夹的角，它与等边三角形的内角互补，即：
  \[
    \theta = \pi - \frac{\pi}{3} = \frac{2\pi}{3}.
  \]
  故选 D。%
}
